\section{十元三平均的十二个精确二次覆盖}\label{sec:ten-cover}
以下给出下降覆盖的全部二次多项式，而不是仅报告程序通过。对表列局部代表 $(X,Y)$，输入尺度及各宏的约分因子按十元构造中的规则决定。表中 $i:(a,b,c)$ 表示第 $i$ 个宏的严格下降条件 $a+bt+ct^2>0$，其中 $t=Y/X$。每行全部开区间均被覆盖，列出的只是未严格覆盖的孤立实点；无穷方向另列。局部相容性通过同一模 $96$ 签名检查，唯一相容例外为 $X=0$，此时已有零坐标，九元网络收尾。
\small\begin{longtable}{r r l p{83mm} p{24mm}}
\toprule 类&个数&代表&$i:(a,b,c)$&未严格覆盖\\\midrule\endhead
1&768&$(0,1)$&1:(-10,-20,-2); 2:(-40,80,0); 3:(0,80,-112); 6:(50,0,-8); 7:(-400,-80,0)&$\infty$\\
2&768&$(1,0)$&1:(-10,-20,-2); 2:(-10,20,0); 3:(0,80,-112); 5:(0,-160,8); 7:(-400,-80,0)&0\\
3&256&$(1,2)$&1:(-10,-20,-2); 2:(-10,20,0); 3:(160,80,-80); 5:(640,-40,130); 7:(-400,-80,0)&$\varnothing$\\
4&384&$(1,3)$&1:(-40,-80,-8); 3:(0,20,-28); 4:(50,20,2); 7:(-100,-20,0)&$\varnothing$\\
5&768&$(1,6)$&1:(-10,-20,-2); 2:(-10,20,0); 3:(0,80,-112); 5:(0,-40,2); 7:(-400,-80,0)&0\\
6&384&$(1,7)$&1:(-40,-80,-8); 3:(0,20,-28); 4:(200,80,8); 7:(-100,-20,0)&$\varnothing$\\
7&256&$(1,8)$&1:(-10,-20,-2); 2:(-10,20,0); 3:(160,80,-80); 5:(2560,-160,520); 7:(-400,-80,0)&$\varnothing$\\
8&128&$(1,11)$&1:(-40,-80,-8); 3:(40,20,-20); 4:(50,20,2); 7:(-100,-20,0)&$\varnothing$\\
9&128&$(1,23)$&1:(-40,-80,-8); 3:(40,20,-20); 4:(200,80,8); 7:(-100,-20,0)&$\varnothing$\\
10&256&$(2,1)$&1:(-10,-20,-2); 2:(-40,80,0); 3:(160,80,-80); 6:(2760,0,480); 7:(-400,-80,0)&$\varnothing$\\
11&768&$(2,3)$&1:(-10,-20,-2); 2:(-40,80,0); 3:(0,80,-112); 6:(200,0,-32); 7:(-400,-80,0)&$\infty$\\
12&256&$(4,5)$&1:(-10,-20,-2); 2:(-40,80,0); 3:(160,80,-80); 6:(690,0,120); 7:(-400,-80,0)&$\varnothing$\\
\bottomrule\end{longtable}\normalsize
每个多项式至多二次。核验一行只须将其全部实根排序，在每个相邻根区间取一个有理数检查符号，再逐根检查等号。根可以写为有理数或有理数加有理倍的平方根；比较根时先分离符号再平方，因而全程不需要浮点容差。二次项系数给无穷方向的符号。由于表列完整的系数及例外集，这一有限证明可独立重做。
